% Generated draft: requires mathematical and source review.
\documentclass[11pt]{article}
\usepackage[margin=25mm]{geometry}
\usepackage{fontspec}
\setmainfont{Latin Modern Roman}
\usepackage{amsmath,amssymb,mathtools,unicode-math}
\setmathfont{Latin Modern Math}
\usepackage{enumitem,needspace,xurl,hyperref}
\setlength{\parindent}{0pt}
\setlength{\parskip}{5pt}
\setlength{\emergencystretch}{3em}
\newcommand{\N}{\mathbb{N}}
\newcommand{\Z}{\mathbb{Z}}
\newcommand{\Q}{\mathbb{Q}}
\newcommand{\R}{\mathbb{R}}
\newcommand{\C}{\mathbb{C}}
\newcommand{\F}{\mathbb{F}}
\newcommand{\A}{\mathbb{A}}
\newcommand{\PP}{\mathbb{P}}
\newcommand{\E}{\mathbb{E}}
\newcommand{\eps}{\varepsilon}
\newcommand{\dd}{\,\mathrm{d}}
\begin{document}
\section*{UnsolvedMath ID: 20002078}
\textbf{Source problem code:} catalogue entry
\subsection{Definitions and mathematical statement}
Question 1.20 has three separate Lagrangian-isotopy problems. Use the canonical cotangent symplectic structure on \(T^*\Sigma_g\) and the standard symplectic form on \(\R^{2n}=\C^n\). Lagrangian isotopy means a smooth family of Lagrangian embeddings; it is not automatically Hamiltonian isotopy.

(a) For \(g>1\), if a closed connected oriented embedded Lagrangian surface \(L\subset T^*\Sigma_g\) satisfies
\[
 [L]=[\Sigma_g]\in H_2(T^*\Sigma_g;\Z),
\]
must it be Lagrangian isotopic to the zero-section? The source reports positive genus-zero and genus-one results of Richard Hind and Alex Ivrii, respectively. The aid supplies a reduction for the higher-genus setting: let \(\lambda\) be the canonical one-form and \(\pi:T^*\Sigma_g\to\Sigma_g\). The degree-one projection and Lagrangian self-intersection identity imply that \(L\) has genus \(g\), and
\[
 (\pi|_L)^*:H^1(\Sigma_g;\R)\longrightarrow H^1(L;\R)
\]
is an isomorphism. Hence choose a closed one-form \(\alpha\) on \(\Sigma_g\) with
\[
 (\pi|_L)^*[\alpha]=-[\lambda|_L].
\]
The fibre translations \(\tau_{s\alpha}(q,p)=(q,p+s\alpha_q)\) are symplectic and carry \(L\) through Lagrangians to an exact Lagrangian, because
\(\tau_\alpha^*\lambda=\lambda+\pi^*\alpha\).
The translation need not be Hamiltonian when \([\alpha]\ne0\). Reaching an exact representative does not prove isotopy to the zero-section; the aid leaves that stronger problem unresolved.

(b) For \(n>2\), is there an embedded Lagrangian torus \(T^n\subset\C^n\) not Lagrangian isotopic to the Clifford torus \(\{|z_1|=\cdots=|z_n|=1\}\)? The source reports a negative answer for \(n=2\) by Ivrii. The aid reports that the corrected 2024 Dimitroglou Rizell--Evans construction gives two monotone Lagrangian \(n\)-tori that are not smoothly isotopic for every even \(n\geq4\). Here monotone means the symplectic-area homomorphism on relative discs is a positive multiple of the Maslov index. If both tori were Lagrangian isotopic to the Clifford torus they would be smoothly isotopic to each other; therefore at least one gives an affirmative answer to (b) in each such even dimension. A distinction only up to Hamiltonian isotopy would not suffice for this question. The record supplies no corresponding conclusion in all odd dimensions \(n>2\).

(c) For \(n>2\), can a Lagrangian submanifold asymptotic to a Lagrangian plane fail to be Lagrangian isotopic to that plane? The source reports a negative answer in dimension \(n=2\) by Eliashberg--Polterovich. The meaning of asymptotic and allowed isotopies is not fixed. The aid's reduction concerns the strong flat-at-infinity version: properly embedded Lagrangian planes \(L\simeq\R^n\) equal to a fixed plane outside a compact set, and isotopies fixed outside one common compact set.

For \(n>1\), any such fixed-end Lagrangian isotopy is, up to reparametrization of \(\R^n\), induced by a compactly supported Hamiltonian isotopy of \(\C^n\). Indeed its deformation one-forms
\[
 \beta_t=F_t^*(\iota_{\partial_tF_t}\omega)
\]
are closed and compactly supported, and \(H_c^1(\R^n;\R)=0\) for \(n>1\). Compactly supported primitives extend to ambient Hamiltonians, with a sign chosen to match the Hamiltonian-vector-field convention. This equates two notions of an existing fixed-end isotopy; it does not produce an isotopy or rule out knots. It also does not treat arbitrary decay-only asymptotic ends. The aid reports (a) and (c) as still unresolved in their natural strong forms; its account does not claim a new current literature assessment.

The original observation that contemporary holomorphic-curve techniques seemed unavailable in higher dimensions belongs to the 2003 workshop context and is not asserted as a present limitation.
\subsection{Short English statement}
The aid reports knotted tori in every even dimension at least four, answering part (b) there. For cotangent surfaces it reduces the homology condition to an exact representative. For flat-ended planes it upgrades any existing fixed-end Lagrangian isotopy to a compactly supported Hamiltonian one. Neither reduction solves parts (a) or (c).
\subsection{Sources}
\begin{itemize}
\item Permanent UnsolvedMath ID \(20002078\), exact public code catalogue entry; local registry \texttt{aim-geometry-notes.json}, \texttt{notes[415]}, canonical tag question.
\item Michael Hutchings, with help from Yasha Eliashberg and John Etnyre, \emph{Holomorphic curves in contact geometry}, AIM workshop problem list, version 15 October 2003, Question 1.20: \url{https://aimath.org/WWN/contactgeom2/contactgeom2.pdf}. The supplied verification report says official HTML and PDF agree; no PDF or HTML was accessed for this draft.
\item Original attributions: Richard Hind for \(g=0\), Alex Ivrii for \(g=1\) and the two-torus isotopy claim, and Eliashberg--Polterovich for the local-plane \(n=2\) result. The aid attributes the corrected 2024 even-dimensional monotone-torus construction to Dimitroglou Rizell--Evans. Full papers and the correction details are not supplied.
\item Research attempt 1, ``Even-dimensional knotted tori and two isotopy reductions,'' labelled partial\_result, partially\_solved, proved and low novelty confidence. Original literature and proposer/year metadata are absent beyond the supplied author/date history.
\end{itemize}
\end{document}
